\documentclass{beamer}

%encoding
%--------------------------------------
\usepackage[T1]{fontenc}
\usepackage[utf8]{inputenc}
%--------------------------------------

%Portuguese-specific commands
%--------------------------------------
\usepackage[portuguese]{babel}
%--------------------------------------

\usepackage{tikz}
\usepackage{pgfplots}
\pgfplotsset{compat = newest}
\usepgfplotslibrary{colormaps}

%Hyphenation rules
%--------------------------------------
\usepackage{hyphenat}
\hyphenation{mate-mática recu-perar}
%--------------------------------------

\usepackage{wrapfig}
\usepackage{graphicx}
\usepackage{biblatex}

% \usepackage[
% backend=biber,
% style=alphabetic,
% sorting=ynt
% ]{biblatex}
\addbibresource{latex-beamer-template.bib}

\graphicspath{ {../../images/} {.images} }

% \usefonttheme{structuresmallcapsserif}
\usetheme{Warsaw}
% \usecolortheme{beetle}
\usecolortheme{seagull}

% \setbeamertemplate{navigation symbols}{}
% \setbeamertemplate{mini frames}{}
\setbeamertemplate{headline}{}
% \renewcommand*{\slideentry}[6]{}
\setbeamertemplate{section in toc}[sections numbered]
\setbeamertemplate{subsection in toc}[subsections numbered]

\setbeamercolor{block body}{bg=yellow!30}
\setbeamercolor{block title}{bg=yellow!40}

\setbeamercolor{block body alerted}{bg=red!10}
\setbeamercolor{block title alerted}{bg=red!20}

\setbeamercolor{block body example}{bg=green!10}
\setbeamercolor{block title example}{bg=green!20}

% Information to be included in the title page:
\title{A Independência do Caminho na Integral do Trabalho}
\subtitle{Trabalho 1 - Grupo 16}

\author [Alex, Caio, Pedro]{
    \small Alex Campbell e Souza - Engenharia de Sistemas \\ 
    Caio Lucas Gomes Silva - Matemática \\ 
    Pedro Mansur Gamarano - Matemática
}

\institute[]{
    \large UFMG \\
    \footnotesize Universidade Federal de Minas Gerais \\
    \small Fundamentos de Eletromagnetismo
}


\date{\today}

\begin{document}

% Slide inicial com título
\frame{\titlepage}

% Índice
%\begin{frame}
%    \frametitle{Index}
\frame{\tableofcontents}
%\end{frame}

% Exibido, no começo de cada seção, o tópico corrente
% \AtBeginSection[]
% {
%   \begin{frame}
%     \frametitle{New Topic}
%     \tableofcontents[currentsection]
%   \end{frame}
% }
\section{sec-1}
\begin{frame}
    \frametitle{Exemplo de Campo Plotado}
    
    \begin{tikzpicture}
        \begin{axis}[
        xmin = -4, xmax = 4,
        ymin = -4, ymax = 4,
        zmin = 0, zmax = 1,
        axis equal image,
        xtick distance = 1,
        ytick distance = 1,
        view = {0}{90},
        scale = 1.25,
        title = {\bf Vector Field $F = [-y,x]$},
        xlabel = {$x$},
        ylabel = {$y$},
        colormap/viridis,
        colorbar,
        colorbar style = {
        ylabel = {Vector Length}
        }
        ]
        \addplot3[
        point meta = {sqrt(x^2+y^2)},
        quiver = {
        u = {-y/sqrt(x^2+y^2)},
        v = {x/sqrt(x^2+y^2)},
        scale arrows = 0.25,
        },
        quiver/colored = {mapped color},
        -stealth,
        domain = -4:4,
        domain y = -4:4,
        ] {0};
        \end{axis}
    \end{tikzpicture}
\end{frame}

% \begin{frame}
%     \frametitle{About me}
%     \begin{wrapfigure}{r}{0.25\textwidth} %this figure will be at the right
%         \centering
%         \caption{My picture. =)}
%         \includegraphics[width=0.25\textwidth]{avatar.jpeg}
%         \label{fig:avatar}
%     \end{wrapfigure}
% 
%     Lorem ipsum dolor sit amet, consectetur adipiscing elit.
% 
%     Curabitur eget tortor at quam ullamcorper pellentesque.
%     Morbi ac enim eu ante viverra elementum id viverra justo.
%     Mauris sagittis mauris sed laoreet maximus. 
%     \cite{dirac}
% 
% %    \begin{figure}[h]
% %        \caption{That's me. =)}
% %        \includegraphics[width=0.5\textwidth, scale=0.25]{avatar.jpeg}
% %    \end{faigure}
% 
%     \listoffigures
% 
% \end{frame}
% 
\section{sec-2}

\begin{frame}
\frametitle{Frame title}

This agent thus does work at a rate $V_1I_1=-{\cal E}_1I_1$. The total work 
    that it does is
    \begin{equation}
        W_1=-\int dtI_1{\cal E}_1=\frac{I_1}{c}\int dt
    \end{equation}
    where $F_1$ is the total change of the flux through the loop. We may 
    also express our result as

\begin{equation}
    a^{2}= b^{2}+c^{2}
\end{equation}

\end{frame}

\section{sec-3}

\begin{frame}
    \frametitle{Highlights}
    
    In this slide, some important text will be
    \alert{highlighted} because it's important.
    Please, don't abuse it.
    
    \begin{block}{Genérico}
    Bloco ``Genérico"
    \end{block}
    
    \begin{alertblock}{Importante}
    Bloco amarelo para ``Atenção".
    \end{alertblock}
    
    \begin{examples}{Boas Práticas}
    Bloco verde para ``Boas Práticas".
    \end{examples}

\end{frame}

\frame{\printbibliography}

\end{document}